\section{Sec.~V: Crossed product by modular group}

Section~IV ended on a genuinely uncomfortable note: type III algebras have no trace, so there is no density
operator and no entropy for them at all — not even the signed, relative-to-a-reference kind that rescued type
II in Sec.~III. If the eventual goal is to compute something like a black-hole entropy using nothing but
operator algebra, this is a real obstruction, not a cosmetic one. The \textbf{crossed product} is the
construction that removes it — not by cheating around the no-trace theorem, but by building a genuinely
\emph{different, larger} algebra, out of the original type III one, that turns out to always be type II (which
does have a trace). The mechanism, remarkably, is completely mechanical and general-purpose, and it will
reappear, essentially unchanged, as the actual physical mechanism behind gravitational entropy starting in
Sec.~IX.

\subsection{Sec.~V, opening: crossed products in general}

Here is the general-purpose construction, stated first without any reference to modular theory at all, because
it's a standard piece of operator-algebra machinery on its own, and it's worth seeing that the specific
application to modular flow (the one this paper actually needs) is just one instance of it. Suppose a group
$G$ acts on a von Neumann algebra $\M$ by unitaries, $\alpha_g(A)=U_gAU_g^\dagger$ (eq.~5.1) — some
one-parameter (or more general) family of symmetries of $\M$. Given the triple $(\M,G,\alpha)$, there is a
standard construction, called the \textbf{crossed product} $\widehat\M\equiv\M\rtimes_\alpha G$, that builds a
\emph{new} von Neumann algebra acting on the enlarged Hilbert space $\widehat\HH=\HH\otimes L^2(G)$ (functions
on $G$, valued in $\HH$ — think of it as attaching an auxiliary quantum system whose configuration space is
$G$ itself). For the case of interest here, $G=\mathbb R$ with generator $K$ and action $\alpha_t(A)=
e^{iKt}Ae^{-iKt}$ (eq.~5.2), so $\widehat\HH=\HH\otimes L^2(\mathbb R)$ — an auxiliary particle on a line.

The specific case this paper needs, and the only one used from here on, takes $K=-\log\Delta_\Psi$ — the
\emph{modular} generator of $\M$ itself, for some cyclic-separating $\ket\Psi$ — so that eq.~5.2 is exactly
the modular flow $\sigma_t$ already built in Sec.~IV. The headline result, proved step by step below: for a
type III algebra $\M$, the resulting crossed product $\widehat\M$ is \textbf{always type II} — regardless of
which type III subtype $\M$ started as. This immediately buys back everything Sec.~III's machinery (density
operators, entropy) needs, applied now to $\widehat\M$ instead of $\M$ itself. And — a fact worth flagging
immediately, even though it's only fully justified at the end of Sec.~V.B — $\widehat\M$ turns out to depend
only on the algebra $\M$, not on which reference state $\ket\Psi$ was used to build the modular flow that
crossed it — so this is really a construction \emph{intrinsic} to $\M$, not an artifact of an arbitrary
choice. This is also, as flagged already back in Sec.~I, exactly the mechanism used in Sec.~IX to explain
black-hole and de~Sitter entropy: there, $\hat q$ below will literally be a physical observer's clock Hamiltonian.

\subsection{Sec.~V.A: construction of \texorpdfstring{$\widehat\M$}{\^M}}

Attach a genuine one-dimensional quantum system — a particle on a line, with position $\hat q$ and momentum
$\hat p$ satisfying the ordinary canonical commutation relation $[\hat q,\hat p]=i$ (eq.~5.3) — to the original
system, giving the enlarged Hilbert space $\HH\otimes L^2(\mathbb R)$. Define the single operator
\[
C \equiv K+\hat q
\]
(eq.~5.4), and let $\widehat\M$ be the subalgebra of $\M\otimes B(L^2(\mathbb R))$ consisting of everything
that commutes with $C$: $a\in\widehat\M$ iff $[a,C]=0$. It's worth pausing on why this particular condition is
the physically right one to impose, since it isn't obvious on sight. $C$ generates \emph{simultaneous}
translations: shifting $\hat q$ by some amount while also flowing modular time by the same amount (since $K$
generates the modular flow and $\hat q$ generates ordinary translations on $L^2(\mathbb R)$, $e^{iCs}$
translates both at once). Demanding invariance under $C$ is exactly demanding that only the \emph{relative}
reading of ``clock minus modular time'' matters, not either one separately — mathematically identical to the
familiar physics idea that only relative position matters, never an absolute coordinate origin. Physically,
in the applications of Sec.~IX, $\hat q$ will be an observer's own energy and $\hat p$ their proper time, and
this invariance condition becomes exactly the statement that physical observables must be
diffeomorphism-invariant (unable to depend on an arbitrary choice of where to put a clock's zero).

Two families of operators manifestly satisfy $[a,C]=0$, and it can be shown they generate the whole of
$\widehat\M$. First, $\hat q$ itself trivially commutes with $C=K+\hat q$ (since $[\hat q,\hat q]=0$ and
$\hat q$ acts on a different factor than $K$), so any (Weyl-form, since $\hat q$ is unbounded) function of it,
$e^{-i\hat qs}$, belongs to $\widehat\M$. Second — and this is the nontrivial part — ordinary elements $A\in
\M$ do \emph{not} commute with $C$ on their own (since $[A,K]\ne0$ in general — that's exactly what it means
for $K$ to generate a nontrivial flow on $A$), but a specific \emph{dressed} version of them does:
\[
[e^{iK\hat p}Ae^{-iK\hat p},\,C] = 0, \qquad A\in\M
\]
(eq.~5.5). Let us prove this commutator identity directly using ordinary canonical commutation relations. Let $\widehat A \equiv e^{iK\hat p}(A\otimes\id)e^{-iK\hat p}$. Using $[\hat q,\hat p]=i$, the position operator acts as a derivative in momentum space: $[\hat q, f(\hat p)] = i f'(\hat p)$. Applying this to the clock translation operator:
\[
[\hat q, \, e^{iK\hat p}] = i(iK)e^{iK\hat p} = -K e^{iK\hat p}, \qquad [\hat q, \, e^{-iK\hat p}] = i(-iK)e^{-iK\hat p} = K e^{-iK\hat p} .
\]
Now compute the commutator $[\hat q, \widehat A]$ using the Leibniz product rule:
\begin{align*}
[\hat q, \, \widehat A] &= [\hat q, \, e^{iK\hat p}](A\otimes\id)e^{-iK\hat p} + e^{iK\hat p}(A\otimes\id)[\hat q, \, e^{-iK\hat p}] \\
&= -K e^{iK\hat p}(A\otimes\id)e^{-iK\hat p} + e^{iK\hat p}(A\otimes\id) e^{-iK\hat p} K \\
&= -K \widehat A + \widehat A K = -[K, \, \widehat A] .
\end{align*}
Therefore, evaluating the commutator with the total constraint $C = K + \hat q$:
\[
[C, \, \widehat A] = [K + \hat q, \, \widehat A] = [K, \, \widehat A] + [\hat q, \, \widehat A] = [K, \, \widehat A] - [K, \, \widehat A] = 0 !
\]
The modular-flow dressing by the clock's momentum $\hat p$ exactly cancels the non-commutativity with $K$, ensuring that $\widehat A$ is strictly invariant under $C$.
So
\[
\widehat\M = \big\{e^{iK\hat p}Ae^{-iK\hat p},\ e^{-i\hat qs} \ \big|\ A\in\M,\ s\in\mathbb R\big\}''
\]
(eq.~5.6), and a general element has the schematic form $\widehat A=\int ds\,A(\hat p;s)\,e^{-i\hat qs}$
(eq.~5.7), with $A(\hat p;s)\equiv e^{iK\hat p}A(s)e^{-iK\hat p}$ built from an operator-valued function
$A(s)\in\M$ — since $\hat p$ is now an operator, $A(\hat p;s)$ is a genuine operator-valued ``function of an
operator,'' living in a kind of quantum spacetime.

There's a second, equivalent way to present the same algebra, obtained by a unitary change of frame
$U=e^{-iK\hat p}$ (which conjugates $\widehat\M\to U\widehat\M U^\dagger$, $C\to UCU^\dagger$), and this
alternative form turns out to be more convenient for essentially everything that follows:
\[
\widehat\M = \big\{A,\ e^{i(K-\hat q)s} \ \big|\ A\in\M,\ s\in\mathbb R\big\}'', \qquad
\widehat A = \int ds\, A(s)\,e^{is(K-\hat q)}, \ A(s)\in\M
\]
(eqs.~5.10, 5.12) — in this frame, ordinary (undressed) elements of $\M$ sit directly inside $\widehat\M$,
and all the dressing has been absorbed into the clock-dependent generator $K-\hat q$ instead.

\subsection{Sec.~V.B: \texorpdfstring{$\widehat\M$}{\^M} is always type II}

\subsubsection*{Finding the modular operator of $\widehat\M$}

Now specialize to $K=-\log\Delta_\Psi$, the genuine modular generator of $\M$, and take the reference vector
\[
\ket{\widehat\Psi} = \ket\Psi\otimes\ket{p=0}
\]
(eq.~5.13, the clock prepared in a definite-momentum, i.e.\ completely spread out in position, state — not
literally normalizable, a technical point flagged in the paper's footnote~30 but not one that affects any
conclusion below). It can be shown $\ket{\widehat\Psi}$ is cyclic and separating for $\widehat\M$, so
Tomita--Takesaki theory applies, and the modular operator $\widehat\Delta$ can be found by directly solving
the defining KMS relation, $\braket{\widehat\Psi|\widehat A\widehat B|\widehat\Psi}=\braket{\widehat\Psi|
\widehat B\,\widehat\Delta\,\widehat A|\widehat\Psi}$ (eq.~5.14), for $\widehat\Delta$ (a computation carried
out in the paper's Appendix~B, not reproduced here). The answer, remarkably simple given how much machinery
went into $\widehat\M$'s construction, is:
\[
\widehat\Delta = \Delta_\Psi
\]
(eq.~5.15) — \textbf{the new algebra's modular operator, in this particular reference state, is exactly the
same operator as the original algebra's modular operator.} Nothing new needed to be invented; the crossed
product inherited its modular structure wholesale.

\subsubsection*{Why this forces type II}

Here is the key algebraic manipulation, and it's short enough to walk through completely. Using
$\widehat\Delta=e^{-K}$ directly:
\[
\widehat\Delta = e^{-K} = e^{-(K-\hat q)}e^{-\hat q} \equiv \rho\rho'^{-1}, \qquad
\rho\equiv e^{-(K-\hat q)}, \quad \rho'\equiv e^{\hat q}
\]
(eq.~5.16), so
\[
\widehat\Delta^{-is} = e^{i(K-\hat q)s}\,e^{i\hat qs}
\]
(eq.~5.17). Now look at what the two factors on the right actually are, comparing against the two presentations
of $\widehat\M$ and $\widehat\M'$ (eqs.~5.10--5.11) given above: $e^{i(K-\hat q)s}$ is manifestly a unitary
sitting inside $\widehat\M$ itself (it's literally one of the generators listed in eq.~5.10), and $e^{i\hat
qs}$ is manifestly a unitary sitting inside $\widehat\M'$ (the commutant, built from eq.~5.11's generators).
\textbf{So $\widehat\Delta^{-is}$ factors as a product of a unitary in $\widehat\M$ times a unitary in
$\widehat\M'$} — which is exactly the criterion (eq.~4.22 of Sec.~IV.B) for the modular flow to be an
\emph{inner} automorphism of $\widehat\M$. But Sec.~IV.B's criterion (eq.~4.21) says modular flow is inner for
\emph{all} $s$ if and only if the algebra is type I or type II — never type III. \textbf{So $\widehat\M$
cannot be type III}, whatever $\M$ itself was. And since $\widehat\HH=\HH\otimes L^2(\mathbb R)$ still cannot
be factorized with respect to $\widehat\M$ (nothing about attaching a clock and imposing an invariance
condition magically produced a tensor factorization — the underlying obstruction from $\M$ is still there),
$\widehat\M$ cannot be type I either. \textbf{By elimination, $\widehat\M$ is type II.} This is worth
appreciating as a genuinely clean piece of reasoning: no explicit trace needed to be constructed to reach this
conclusion — it followed purely from identifying the modular operator and checking the inner-automorphism
criterion already established in Sec.~IV.

\subsubsection*{Building the trace explicitly, and checking cyclicity by hand}

Having established \emph{that} a trace exists, it's worth actually writing one down and checking it really
works, rather than just trusting the general argument. Define
\[
\tr\widehat A \equiv \braket{\widehat\Psi|\widehat A\,\rho^{-1}|\widehat\Psi}
\]
(eq.~5.20, using $\rho=e^{-(K-\hat q)}$ from eq.~5.16 — motivated by the standard fact that if a trace exists,
expectation values in \emph{any} state should look like $\braket{A}=\tr(A\rho_{\text{state}})$ for some density
operator $\rho_{\text{state}}$; solving for what $\rho_{\text{state}}$ would need to be, given that
$\ket{\widehat\Psi}$ is itself the reference state used to build $\Delta_\Psi=\widehat\Delta$, directly
motivates dividing by $\rho$). Cyclicity, $\tr(\widehat A\widehat B)=\tr(\widehat B\widehat A)$, is checked
directly using nothing but the KMS relation (eq.~5.14) that defined $\widehat\Delta$ in the first place:
\[
\tr(\widehat A\widehat B) = \braket{\widehat\Psi|\widehat A\widehat B\rho^{-1}|\widehat\Psi}
= \braket{\widehat B\rho^{-1}\widehat\Delta\,\widehat A} = \braket{\widehat B\rho'^{-1}\widehat A}
= \braket{\widehat B\widehat A\rho'^{-1}} = \tr(\widehat B\widehat A)
\]
(eq.~5.22, suppressing the common $\braket{\widehat\Psi|\cdots|\widehat\Psi}$ for brevity in each term): the
first step is just the definition of $\tr$; the second is the KMS relation (eq.~5.14) applied with the roles
of $A\rho^{-1}$ and $B$ swapped; the third uses $\rho^{-1}\widehat\Delta=\rho^{-1}\rho\rho'^{-1}=\rho'^{-1}$
directly from eq.~5.16's factorization; and the fourth uses that $\rho'^{-1}=e^{-\hat q}\in\widehat\M'$
(established just above) commutes with $\widehat A\in\widehat\M$ by definition of the commutant. Every step is
either a definition or something already established — no new assumption anywhere. This is worth comparing to
the very similar, but much more concrete, cyclicity check already carried out by hand in Sec.~III.C.1
(eqs.~3.12--3.14 of that companion section, for the finite Bell-pair-chain trace): the mechanism is
structurally the same idea (a specific reference state's own consistency conditions force cyclicity), just
now stated in the fully general, type-III-compatible language of modular theory rather than worked out on
explicit small matrices.

\subsubsection*{Type $\mathrm{II}_\infty$, or type $\mathrm{II}_1$ if you restrict the clock's energy}

Compute $\tr(\id)$ directly from eq.~5.21 (an explicit rewriting of eq.~5.20 for the identity element):
\[
\tr(\id) = \int_{-\infty}^{\infty} dq\, e^{-q} \braket{\Psi|\id|\Psi} = \int_{-\infty}^\infty dq\,e^{-q} = \infty
\]
(eq.~5.25) — divergent, so by the Sec.~II.C classification, $\widehat\M$ is (so far) type $\mathrm{II}_\infty$.
But here is a genuinely simple, physically motivated fix, worth seeing in full because it's exactly the
mechanism reused in Sec.~IX: restrict the clock's spectrum to be bounded below, $q\ge0$ — physically, the
statement that a real observer's energy cannot be negative, certainly a reasonable restriction on any
physically sensible clock. Concretely, insert the projector $\Pi=\theta(\hat q)$ (the step function, $1$ for
$q\ge0$ and $0$ for $q<0$) and define $\widehat\M_+\equiv\Pi\widehat\M\Pi$, acting only on the restricted
Hilbert space $\HH\otimes L^2(\mathbb R_{>0})$. The identity of this smaller algebra is $\Pi$ itself, and
\[
\tr_{\widehat\M_+}(\id) = \tr_{\widehat\M}(\Pi) = \int_{-\infty}^\infty dq\,e^{-q}\theta(q) = \int_0^\infty
dq\,e^{-q} = 1
\]
(eq.~5.26) — finite, and in fact already normalized to $1$ with no further rescaling needed. \textbf{So
restricting a physical observer's clock to have positive energy is exactly the algebraic operation that turns
type $\mathrm{II}_\infty$ into type $\mathrm{II}_1$} — a purely mathematical fact about which subtype of type
II you land in, but one with an immediate and physically loaded reading: whether the observer crossing your
algebra has a bounded-below energy spectrum directly decides which flavor of type II entropy you'll be
computing. This exact fork reappears, with real physical stakes attached, when black hole entropy
(type $\mathrm{II}_\infty$, Sec.~IX.B) is contrasted with de~Sitter entropy (type $\mathrm{II}_1$, Sec.~IX.C).

\subsubsection*{The crossed product doesn't actually depend on which reference state you started with}

One loose end, promised at the start of this section: $\widehat\M$ was built using a specific cyclic-separating
$\ket\Psi$ to define the modular generator $K$ that got crossed with — does a different choice $\ket\Phi$ give
a genuinely different algebra? It can be shown it does not, only a \emph{unitarily equivalent} one,
\[
\widehat\M_\Phi = u'_{\Phi\Psi}(\hat p)\,\widehat\M_\Psi\,u_{\Phi\Psi}'^\dagger(\hat p)
\]
(eq.~5.27), where $u'_{\Phi\Psi}$ is exactly the intertwining unitary from Sec.~IV.F (eq.~4.114), with the flow
parameter identified with the clock's momentum operator $\hat p$. (The short computation establishing this,
eq.~5.29, is a direct if slightly involved manipulation using the identities from Sec.~IV.F and is not
reproduced symbol-by-symbol here — the content that matters is the conclusion: $\widehat\M$ is, up to unitary
equivalence, an \emph{intrinsic} invariant of $\M$ alone, exactly as claimed at the start of this section.)

\subsection{Sec.~V.C: density operator for \texorpdfstring{$\widehat\M$}{\^M} in a general semiclassical state}

With a genuine trace in hand, Sec.~III's machinery — a density operator defined by
$\tr(A\rho_{\widehat\M})=\braket{\text{state}|A|\text{state}}$ — becomes available for $\widehat\M$. The
paper works this out for a
class of physically motivated states of the form $\ket{\widehat\Phi}=\ket\Phi\otimes\ket g$ (eq.~5.30), where
$\ket g$ is some fixed clock wavefunction, normalized ($\int dq\,|g(q)|^2=1$, eq.~5.31) and nonvanishing
everywhere. Solving the defining equation for the resulting density operator is a genuinely involved
computation (eqs.~5.32--5.41 of the paper, using the two-state KMS relation, eq.~4.105, from Sec.~IV.F, twice
over, together with the explicit form $\braket{q|e^{-iK\hat p}|g}=g(q-K)$, eq.~5.37, for how the clock
wavefunction transforms between the two frames of eqs.~5.6 and 5.10) — not walked through symbol by symbol
here, since the manipulation itself is mostly bookkeeping once the KMS relation is trusted, but the outcome is
clean:
\[
\rho_{\widehat\Phi} = 2\pi\, g(\hat q-K)\,e^{\hat q}\,\Delta_{\Phi\Psi}\, g^*(\hat q-K)
\]
(eq.~5.39, in the frame of eqs.~5.6--5.7; the corresponding expression in the other frame, eq.~5.10, is
obtained by conjugating with $U^\dagger$, eq.~5.41) — an explicit operator, built from the clock wavefunction
$g$, the (already familiar) $e^{\hat q}$ factor from the trace's own normalization, and the relative modular
operator $\Delta_{\Phi\Psi}$ from Sec.~IV.F comparing the state of interest $\Phi$ to the reference $\Psi$
used to build the crossed product in the first place.

\subsection{Sec.~V.D: entanglement entropy for \texorpdfstring{$\widehat\M$}{\^M} in a general semiclassical state}

\subsubsection*{The final formula, and why it has exactly the shape it does}

Plugging $\rho_{\widehat\Phi}$ into the ordinary entropy formula, $S_{\widehat\M}=-\tr(\rho_{\widehat\Phi}\log
\rho_{\widehat\Phi})$ (eq.~5.42), and specializing to a \textbf{semiclassical} clock wavefunction — one that
varies slowly, $g'(q)\propto\epsilon$ for some small parameter $\epsilon$, and working to leading order in
$\epsilon$ — the paper carries out an expansion (eqs.~5.43--5.48: expand $-\log\rho_{\widehat\Phi}$ using
$K_{\Phi\Psi}=K_\Phi+K-K_{\Psi\Phi}$ from Sec.~IV.F's eq.~4.116, then evaluate the resulting expectation value
term by term, using $K_\Phi\ket\Phi=0$ from the ordinary modular-operator property eq.~4.9, and recognizing
the surviving piece directly as the relative entropy definition, eq.~4.117, of Sec.~IV.F) that lands on
\[
S_{\widehat\M}(\widehat\Phi) = -S_\M(\Phi\|\Psi) - \bar q + S_o, \qquad
\bar q = \int dq\,q\,|g(q)|^2, \qquad S_o = -\int dq\,|g(q)|^2\log|g(q)|^2
\]
(eqs.~5.49--5.50): $\bar q$ is simply the mean clock reading under the probability distribution $|g(q)|^2$, and
$S_o$ is the ordinary (differential) Shannon entropy of that same distribution — both depending only on the
shape of the clock wavefunction $g$, not on the state $\Phi$ being described. \textbf{The entropy of the type
II crossed-product algebra, up to state-independent constants fixed entirely by the clock, is exactly minus
the type III relative entropy of the original algebra $\M$.} This is worth appreciating as a genuinely
satisfying closing of the loop opened at the very start of Sec.~IV: relative entropy, introduced there as ``the
one thing that survives type III when ordinary entropy doesn't,'' now turns out to literally \emph{be} an
ordinary, honest entropy — just of a different, larger algebra, built by attaching a clock.

\begin{quote}
\textit{A concrete check of the general shape of this formula, worked with an explicit Gaussian clock.} It's
worth seeing the qualitative content of eq.~5.49 confirmed on a fully explicit, computable example, even
though the actual derivation above is more careful about tracking exactly which piece is the relative-entropy
term. Model the leading-order, unnormalized weight the state assigns to clock reading $q$ as
$\rho(q)\propto g(q)^2\,e^{q}$ (the $e^q$ factor tracking the $\rho^{-1}=e^{K-\hat q}$ dependence built into the
trace's own definition, eq.~5.20, once the $\M$-sector operator content has been integrated out against a
fixed state and only the clock dependence is left), for a Gaussian clock wavefunction of width $\sigma$
centered at $\bar q_0$, $g(q)^2=\frac{1}{\sqrt{2\pi}\sigma}\exp\!\big(-(q-\bar q_0)^2/2\sigma^2\big)$. Carrying
out the resulting Gaussian integral symbolically: the normalization is $Z=\int dq\,g(q)^2e^{q}=\exp(\bar
q_0+\sigma^2/2)$, and the resulting (still Gaussian, same width $\sigma$) tilted distribution
$\rho(q)/Z$ has its mean shifted up by exactly one unit of variance, $\braket{q}=\bar q_0+\sigma^2$ — the
$e^q$ factor systematically biases the observed clock reading upward, an effect you can see directly in this
toy model rather than having to trust it abstractly. Its (differential) entropy is the standard Gaussian
formula $\tfrac12\log(2\pi e\sigma^2)$, so
\[
-\braket{q} + \big(\text{entropy}\big) = -\bar q_0 - \sigma^2 + \tfrac12\log(2\pi e\sigma^2) ,
\]
computed exactly for this Gaussian example: an expression of the form $-\bar q_0$ plus clock-shape-only terms,
matching the general shape of eq.~5.49 (with the leftover $-\sigma^2+\tfrac12\log(2\pi e\sigma^2)$ terms
playing the role of the general formula's $S_o$, in this simplified toy where the relative-entropy piece
has been set aside by construction). This is offered as a sanity check on the \emph{shape} of the general
result, not as a substitute for the paper's own, more careful derivation above, which correctly separates out
the $\M$-sector relative-entropy contribution that this simplified toy model doesn't track.
\end{quote}

\bigskip
\noindent This closes Sec.~V, and with it, the entire first half of the paper (Secs.~II--V): a type III
algebra, with no trace and no entropy of its own, can always be crossed by its own modular group to produce a
type II algebra that does have both — and that type II algebra's entropy is nothing but the original algebra's
relative entropy, in disguise. Section~VI turns to the payoff this machinery was built for: applying every
tool developed so far — the type classification, modular theory, and now the crossed product — to the
large-$N$ limit of the AdS/CFT correspondence itself.
